\newpage
\appendix
\section{替代的 CT18Z 全球拟合
\label{sec:AppendixCT18Z}
}

\begin{figure}[b]
	\includegraphics[width=0.49\textwidth]{images/sqrt2chi2_ct18ZnnloCount.pdf}
	\caption{与图~\ref{fig:sn_ct18} 类似，全部 CT18Z 数据集的有效高斯变量（$S_E$）分布。两个方块与两颗星分别表示 NuTeV 双缪子与 CCFR 双缪子实验的 $S_E$ 值。
\label{fig:sn_ct18z}}
\end{figure}

本附录描述一系列通向 CT18Z PDF 的拟合；CT18Z 为本分析的主要结果 CT18 NNLO 提供了截然不同的备选。
虽然 CT18Z NNLO 在描述 CTEQ-TEA 数据上达到相当的成功水平——$\chi^2/N_\mathit{pt}\! =\! 1.19$，对比 CT18 NNLO 拟合的 $\chi^2/N_\mathit{pt}\! =\! 1.17$——具体数据集的符合质量发生了多处变化。这部分可通过比较 CT18Z 的 $S_E$ 值分布（图~\ref{fig:sn_ct18z}）与我们为 CT18 给出的结果（图~\ref{fig:sn_ct18}）看到。
在 CT18Z 中纳入 ATLAS 7 TeV $W/Z$ 数据（即 ATL7ZW 数据，CTEQ 实验 ID=248）使许多实验的 $S_E$（或 $\chi^2_E/N_{pt,E}$）向上移动，尤其是 NuTeV（Exp.~IDs=124、125）与 CCFR（Exp.~IDs=126、127）的双缪子数据，表明 ATL7ZW 数据与其他数据集存在一定分歧。

Sec.~\ref{sec:alt} 已指出，我们发布两个中间拟合：CT18A（仅加入 ATL7ZW 数据）与 CT18X（DIS 截面采用特殊挑选的因子化标度）。相比 CT18A 与 X，CT18Z PDF 使 PDF 及其矩、部分子光度和标准烛光预言相对 CT18 产生最大变化：见图s.~\ref{fig:PDFbands1}、\ref{fig:DBandSBbands}、\ref{fig:DBandSBbands2}、\ref{fig:lumia}、\ref{fig:corr_ellipse1} 与 \ref{fig:corr_ellipse2}。


下面更详细地考察从 CT18 演化到 CT18Z NNLO 的各项修改之积累。Sec.~\ref{sec:CT18Z_changes} 回顾 ATL7ZW 与 CDHSW 实验的作用。Sec.~\ref{sec:CT18ZvsCT18} 总结 CT18、A、X、Z PDF 的关键差异；CT18A 与 X 的 NNLO、NLO PDF 误差带图收录于补充材料。
Sec.~\ref{sec:CT18Z_qual} 探讨与 ATL7ZW 数据的符合情况。
Sec.~\ref{sec:LMCT18Z} 考察 $\chi^2$ 扫描，提取 CT18Z 全球拟合内 PDF 约束重新分配的详细信息以及 CT18Z 对 $\alpha_s(M_Z)$ 与 $m_c$ 的预言。

本文给出的物理结论已用若干独立技术核实。最初，通过应用快速 Hessian 技术 \texttt{PDFSense}/$L_2$
敏感度~\cite{Wang:2018heo,Hobbs:2019gob}
与 \texttt{ePump}~\cite{Hou:2019gfw}，从基于先前 \CTHERAII{} NNLO
PDF~\cite{Hou:2016nqm} 的理论预言出发，获得数据集对 PDF 可能影响的预估。在这一阶段我们发现，剖析 ATL7ZW 数据时，\texttt{ePump} 与 \texttt{xFitter} 程序给出非常不同的结果。该分歧在 App.~\ref{sec:Appendix4xFitter} 中处理。
作为拟合本身的一部分，我们多次重复某些拟合：或用拉格朗日乘子（LM）把 PDF 约束在特定数值，或变动 ATL7ZW 及其他数据集的统计权重，以按 Collins 与 Pumplin~\cite{Collins:2001es} 的方法考察其相互一致性。所有这些方法给出一个连贯的物理图像，现总结如下。

\subsection{对数据集与理论设定的改动}
\label{sec:CT18Z_changes}

\subsubsection{修改后的数据选择}
先讨论两个数据集描述中出现的一些问题：({\it i}) ATLAS 新近测得的 7 TeV Drell-Yan 数据，涉及单举 $W,Z$ 玻色子产生的快度分布（ATL7ZW，Exp.~ID=248）；({\it ii}) CDHSW 从 $\nu$-Fe 数据提取的 $F^p_2$、$xF^p_3$ DIS 结构函数信息（Exp.~IDs=108、109）。

\begin{figure}[tb]
\includegraphics[width=0.7\textwidth]{images/248_ct18znn_L2_q100_Sf_1.pdf}
\vspace{-2ex}
\caption{
	ATL7ZW 数据对几个单独部分子味 PDF 的 $L_2$ 敏感度。对奇异分布
	$s(x,Q)$ 的拉动尤其大，在 $x\!\sim\! 0.02-0.05$ 达到峰值 $S_{s,\mathit{L2}}\! \sim\! -20$；不过对 $d$-、$\bar{u}$-
	与 $\bar{d}$-夸克 PDF 的相反拉动在同一 $x$ 区域也很显著。
}
\label{fig:L2_248}
\end{figure}


\begin{figure}[htbp]
        \begin{center}
                \includegraphics[width=0.95\textwidth]{images/ifl3_ct18znn_L2_q2_Sf_1.pdf}
        \end{center}
        \vspace{-2ex}
        \caption{
		CT18Z 所拟合数据对 $Q\!=2$ GeV 处依赖 $x$ 的奇异 PDF 的 $L_2$ 敏感度。鉴于 $R_s(x,Q)$ 的拟合行为在很大程度上由 $s(x,Q)$ 驱动，此处所示信息与图~\ref{fig:lm_rsz} 的 LM 扫描互补。
        }
\label{fig:L2sZ}
\end{figure}

{\bf ATLAS 7 TeV 单举 $W/Z$ 产生数据。}
就情形 ({\it i}) 而言，ATL7ZW 数据集被认为提供了关于轻夸克核子海结构的重要信息，特别是偏好比以往以 DIS 数据为主导的全球分析更大的奇异性压低因子 $R_s(x,Q)$（式~(\ref{eq:Rs})）。
敏感度分析表明，尽管 ATL7ZW 测量相对各 PDF 味的关联余弦相当温和（典型地 $|\cos \phi|\! <\! 0.6$），ATL7ZW 数据的高精度仍对所有轻反夸克味造成显著拉动：$\bar u$、$\bar d$，尤其 $\bar s$。例如这些拉动可见于图~\ref{fig:L2_248} 中展示 ATL7ZW 数据对各 CT18Z PDF 味的 $L_2$ 敏感度的图示。图~\ref{fig:L2_248} 中 $s(x,Q)$ 在 $x=0.01-0.1$ 处相应的负 $S_{L2,E}$ 所揭示的对 $s(x,Q)$ 的强正向拉动，必须在全球拟合中被其他实验的反向拉动补偿。例如，可以在图~\ref{fig:L2sZ} 绘出各实验分别的敏感度，比较各实验对 $s(x,Q=2\mbox{ GeV})$ 拉动的 Hessian 估计。显而易见，ID=248 在 $x=0.03$ 处对 $\Delta \chi^2$ 的显著负拉动受到 NuTeV（Exp.~ID=124、125）与 CCFR
（Exp.~ID=126、127）双缪子 SIDIS 正向拉动的对抗，还非常突出地受到 HERAI+II 单举数据（Expt.~ID=160）的对抗。另一方面，固定靶 E866 在 $pp$ 靶上的 Drell-Yan 数据（Expt.~ID=204）在与 ATL7ZW 相同的方向上弱弱拉动，只是位于较小
$x \approx 0.01$ 处。

$L_2$ 敏感度估计对 $\chi^2$ 的贡献；与 ATL7ZW 数据的张力也由其他统计指标揭示，例如实验 $E$ 的有效高斯变量 $S_E$，它量化 $\chi^2_E$ 的变化与其统计不确定度的比较。按这一度量，基于 \texttt{ePump} 的 Hessian 更新研究~\cite{Hou:2019gfw} 发现，把 ATL7ZW 数据以递增统计权重纳入 CT18 拟合会强烈抬高 NuTeV $\nu$ SIDIS（Exp.~ID=125）、E866 $\sigma_{pd}/(2\sigma_{pp})$ Drell-Yan 数据（Exp.~ID=203）与 CMS 7 TeV 电子不对称数据（Exp.~ID=267）的 $S_E$ 值，
参见~\cite{Hou:2019gfw} 图 25。这一变化伴随着 $s$ 夸克 PDF 与 $\bar{d}/\bar{u}$ PDF 比值的修改，参见同文献图 26。最后我们观察到，纳入 ATL7ZW 数据集后，$g(x,Q)$ 在 $x \gtrsim 10^{-2}$ 受到轻度压制。

{\bf CDHSW 数据。}
我们的 LM 扫描（如图~\ref{fig:LMg18} 所给者）与图~\ref{fig:L2glu} 的 $L_2$ 敏感度图揭示：铁靶带电流中微子相互作用深度非弹性散射的 CDHSW 测量
（Exp.~IDs=108~\cite{Berge:1989hr}, 109~\cite{Berge:1989hr}）
经由其截面的 $Q^2$
分布而对 $x > 0.2-0.5$ 的胶子分布敏感。在 $x <0.4$，CDHSW
与 CCFR 测得的结构函数 $F_2$、$xF_3$ 的对数斜率不同，参见文献~\cite{Seligman:1997fe} 图 8.3 与 8.4：CDHSW 结构函数偏好 $x\!\approx\!0.2$ 处更硬的 $g(x,Q\!=\!100\,\mathrm{GeV})$，但在 $x\!\gtrsim\!0.5$ 处胶子更软——与 CDF Run-2 喷注数据（Exp.~ID=504）类似。众所周知，CDHSW 分析可能存在未解决的实验问题~\cite{Barone:1999yv}，且铁靶的核修正可能不可忽略。因此人们会问：若剔除这两个 CDHSW 数据集，PDF 会如何变化？该问题曾在文献~\cite{Hou:2019gfw} 用 \texttt{ePump} 处理过，也通过执行名为
``CT18mCDHSW''的特殊拟合（即 CT18``减去''CDHSW）处理，其中两个 CDHSW 数据集被移出 CT18 全球数据。

补充材料所载 CT18mCDHSW NNLO 拟合的 PDF 误差带图表明：剔除 CDHSW 数据导致 $x=0.1-0.5$ 胶子 PDF 略降，同时胶子 PDF 不确定度略增，且同一 $x$ 区域 $u,d$ PDF 出现补偿性增加。

\subsubsection{依赖 $x_B$ 的标度与修改后的全球拟合}
%
另一个潜在关切是 QCD 截面对重正化与因子化标度的残余依赖；与最新的实验不确定度相比，我们在一些实验中发现其不可忽略，即便使用 QCD 截面的 NNLO 表达式。特别地，按表~\ref{tab:AXZ} 在依赖于 Bjorken $x_B$ 的精心选择的因子化标度 $\mu_{F,x}$ 下评估 NNLO DIS 截面，可适度改善与 HERA DIS 截面的符合：
相应 $\chi^2$ 改善 40 个单位（$N_{pt}=1120$ 个数据点上 $\chi^2(\mbox{CT18})=1408$、$\chi^2(\mbox{CT18Z})=1378$），若如右图~\ref{fig:saturation} 把单举 HERA 数据的统计权重增至 10，还可再改善 30 个单位。在此标度选择下，动量分数 $x$ 低于 0.05 处的胶子与奇异性 PDF 也更大，参见图~\ref{fig:saturation} 左图。同时为保持求和规则，$g(x,Q)$ 与 $s(x,Q)$ 从 $x=0.2$ 到 0.4-0.5 被减小。

于是，全球拟合框架中的三项修改——DIS 使用 $\mu_{F,x}$ 标度、剔除 CDHSW 数据集、加入 ATLAS7ZW 数据——加总起来压制了 $0.005 \lesssim x \lesssim 0.3$（即 LHC 经 $gg$ 融合产生希格斯的绝大部分相关 $x$ 区间）的胶子 PDF。它们的组合还在全 $x$ 显著增大 $s(x,Q)$。

如 Sec.~\ref{sec:alt} 所述，这一组合被用于生成 CT18Z NNLO PDF。中间 PDF 组 CT18A 与 CT18X 分别只实现 ATL7ZW 数据的 $66 < Q < 116$ GeV 区域（34
个数据点）或只实现 DIS 标度 $\mu_{F,x}$，见表~\ref{tab:AXZ}。[低亮度
（$35\, {\rm pb}^{-1}$）的 ATLAS 7 TeV $W^\pm,Z$
截面数据样本（Exp.~ID=268）被移出 CT18A 与 Z 拟合以避免重复计数。]

\subsection{四个 PDF 组合之间的比较}
\label{sec:CT18ZvsCT18}

补充材料包含一系列比较 CT18、A、X、Z 各 PDF 味 NNLO 与 NLO 不确定度带的图形。这些比较的主要特征可以提炼如下。

\begin{enumerate}
\item {\bf $g(x,Q)$：}一方面比较 CT18、CT18A 与 CT18Z，另一方面比较 CT18、CT18X 与 CT18Z 的 PDF 不确定度带，显然 CT18Z 偏离 CT18 的主体源于修改的 DIS 标度选择 $\mu_{F,x}$，剔除 CDHSW 与纳入 ATL7ZW 测量的相互作用带来的变化较弱。NLO 下偏离 CT18 一般更小，只有极低 $x$ 区域例外——那里 CT18Z 与 CT18X 相对 CT18 的比值大 $\sim\! 50\%$。
%
\item {\bf $d(x,Q)$：}与 $g(x,Q)$ 相反，Drell-Yan 过程对 $d(x,Q)$ 的敏感度放大了该味上 CT18 与 CT18A NNLO 结果的差异，表现为 $x\!\sim\!10^{-3}$ 附近中心 CT18A 分布相对 CT18 约 $\sim\!1\%$ 的轻度压制，伴随 PDF 不确定度降低百分之几。此区域内的 CT18Z 拟合被向相反方向拉动：$x\!<\! 0.1$ 处增强约 $\lesssim\!2\%$，不确定度相当。在高 $x$ 区域
$x \gtrsim 0.2$，仅在 CT18A NNLO 拟合 ATL7ZW 数据的效果归结为 $d(x,Q)$ 不确定度的小幅移动。相比之下，CT18Z 在更高 $x$ 处相对 CT18 受抑。
与 NNLO 相比，NLO 下拟合 $d(x,Q)$ 的定性影响较弱，主要见于较低 $x\!\lesssim\!
10^{-3}$。
%
\item {\bf $u(x,Q)$ 与 $\bar u(x,Q)$：}这里显著的特性是 CT18Z NNLO 偏离 CT18 NNLO 的程度有多大程度由 CT18X 中的修改驱动：CT18Z 与 X 的中心 PDF 和不确定度高度一致即为明证。一旦在 ATL7ZW 测量之上再把 CT18X 发现的 DIS 标度选择实现于 CT18Z，CT18A NNLO 中 $\bar{u}(x,Q)$ 的弱压制几乎完全消失。
\item {\bf $s(x,Q)$ 与 $R_s(x,Q)$：}前已指出，在 CT18A/Z NNLO 拟合引入 ATL7ZW 数据使 $s(x,Q)$ 明显高于 CT18，而 CT18X 的标度选择在 $0.01 \lesssim x \lesssim 0.5$ 轻度压制 $s(x,Q)$、在 $x \lesssim 0.01$ 与 $x \gtrsim 0.5$ 增强。奇异性压低比 $R_s(x,Q)$ 主要由 $s(x,Q)$ 本身观察到的模式决定。
\end{enumerate}


\subsection{细看 ATLAS 7 TeV $W/Z$ 数据的描述}
\label{sec:CT18Z_qual}


\subsubsection{CT18A(Z) 中对 ATL7ZW 数据的拟合优度}
\label{sec:ATL7ZWchi2}

\begin{table}[tb]
\caption{不同组的 QCD 分析对 ATLAS 7 TeV $W/Z$ 数据 $\chi^{2}$ 值的比较。对 CT18A/Z PDF，我们给出通过对 ATLAS $W/Z$ 数据集的 $\chi^2$ 权重作拉格朗日乘子扫描得到的 68\% C.L. 不确定度，见图~\ref{fig:lm_wht}（左）。}
\label{tab:CMN}
\begin{tabular}{c|c|c|c|c}
\hline\hline
PDF & $N_{\mathit{pt, E}}$ & $\chi^{2}_E$ & $\chi^{2}_E/N_{\mathit{pt}, E}$ & Ref. \\
\hline
CT18A &  34  & $88^{+68}_{-28}$  & $2.58^{+2.01}_{-0.84}$ & \\
CT18Z & 34 & $89^{+65}_{-31}$ & $2.61^{+1.90}_{-0.91}$  & \\
\hline
ATLAS-epWZ16 &  61 & 108    & 1.77 & \cite{Aaboud:2016btc}\\
MMHT (2019)  &  61 & 106.8  & 1.76 &  \cite{Thorne:2019mpt}\\
NNPDF3.1     &  34 & 73     & 2.14 & \cite{Ball:2017nwa} \\
\hline\hline
\end{tabular}
\end{table}

现在转向 ATL7ZW 测量的总体描述——正是它在 CT18A/Z NNLO 中引起核子奇异性的强烈修改。我们发现，这些数据一经实际拟合，其 $\chi^2$ 显著改善：CT18 NNLO 给出非常大的 $\chi^2_E/N_\mathit{pt,E}\!=\! 8.4$（$S_E\! =\! 13.7$），而在 CT18Z NNLO 大幅降至 $\chi^2_E/N_\mathit{pt,E}\!=\! 2.6$（$S_E\! =\! 4.8$），见表~\ref{tab:EXP_2}。注意 CT18A NNLO（其设定与 CT18 NNLO 只差 ATL7ZW 数据的实现）的相应数值仅比 CT18Z 略有改善：$\chi^2_E/N_\mathit{pt,E}\!=\! 2.6$（$S_E\! =\! 4.7$）。还要提醒读者：CT18A/Z 拟合只包含共振区（$66 < Q < 116$ GeV）内的 34 个 ATL7ZW 数据点；已发表集合其余 61 点中的部分精度较低，主要贡献于 $\chi^2_E/N_{pt,E}$ 的降低而非真正改善 PDF 约束，同时还会增强离共振区 PDF 对 NLO EW 修正的依赖。

作为比较，ATLAS 组自身在结合 $W/Z$ 与 HERA I+II 合并数据的 ATLAS-epWZ16 拟合~\cite{Aaboud:2016btc}~\cite{Abramowicz:2015mha} 中得到 $\chi^2_E/N_{\mathit{pt,E}}=108/61$。对 CT14 PDF，他们在 \texttt{xFitter} 中用 $W/Z$ 数据剖析 CT14 PDF 后得到 $\chi^2_E/N_{\mathit{pt,E}}=103/61=1.69$。其拟合的奇异性分数为
$R_s=(s+\bar{s})/(\bar{u}+\bar{d})=1.13\pm0.08$（$x=0.023$），显著大于我们在 CT18A(Z) 拟合所得。

当在 CT 分析中只纳入 ATL7ZW 与 HERA I+II 数据时，我们确认了上述发现。事实上，由于 CT 参数化更灵活，若以更高的统计权重纳入 ATL7ZW 数据，我们会得到更好的 $\chi^2/N_{pt,E}\approx 1.4$ 甚至更大的 $R_s$。

\begin{figure}
    \centering
    \includegraphics[width=0.7\textwidth]{images/CT18_CT18A_NNPDF31.pdf}
    \caption{
        $Q\!=\!100$ GeV 处拟合奇异 PDF $s(x,Q)$ 的比较：
		CT18（紫色实线）与 CT18A（灰色短划线）
		及 NNPDF3.1（品红长划线）。}
    \label{fig:CT18A2NNPDF31}
\end{figure}

然而，出于如下理由我们认为该结果有问题：
\begin{enumerate}
\item 上述 \texttt{xFitter} 剖析分析强烈淡化了与 ATL7ZW 数据存在张力的实验，解释见 App.~\ref{sec:Appendix4xFitter}；
\item 像 Sec.~\ref{sec:LMCT18Z} 所给那样的 LM 扫描显示，与 HERAPDF 拟合相比灵活性更强的 CT 拟合，当被迫让 $R_s(x,Q)$ 在 ATL7ZW 数据偏好的 $x=0.01-0.1$ 接近 1 时会变得不稳定或出现多个极小；这种不稳定或许反映数据区分 $\bar s$、$\bar d$ 与 $\bar u$ 贡献的能力仍然薄弱。
\end{enumerate}

其他 PDF 拟合组也研究过 ATL7ZW 数据。表~\ref{tab:CMN} 总结 CT18A、Z、ATLAS-epWZ16、MMHT（2019）与 NNPDF3.1 NNLO 拟合中 ATL7ZW 数据的 $\chi^{2}_E/N_{pt,E}$ 值。对 CT18A/Z 拟合，我们引用 $\chi^2_{{\rm \small ATL7ZW}}$ 的 68\% C.L. 不确定度，它们来自 Sec.~\ref{sec:ATL7ZWchi2wt} 中对 $\chi^2_{{\rm \small ATL7ZW}}$ 权重的 LM 扫描（图~\ref{fig:lm_wht} 左面板）。不确定度等于 $\chi^2_{{\rm \small ATL7ZW}}$ 与最优全局拟合处取值之差，条件是总 $\chi^2$（图~\ref{fig:lm_wht} 左面板实黑线）增加一个标准差（36 个 $\chi^2$ 单位）。

表~\ref{tab:CMN} 所引各 $\chi^2_{{\rm \small ATL7ZW}}$ 值都很好地落在 CT18A/Z 的 68\% 不确定度之内；名义 CT18A/Z 的 $\chi^2_{\rm \small ATL7ZW}$ 偏高但不显著。$\chi^2_{\rm \small ATL7ZW}$ 的差异主要可追溯到中等 $x$ 处 $s(x,Q)$ 的大小：CT18A/Z NNLO 中它在不确定度之内低于某些其他拟合。

ABM 分析~\cite{Alekhin:2017olj} 强调过 ATL7ZW、NuTeV 与 NOMAD 数据集之间的张力，以及偏好 $R_s(x,Q)$ 增强对奇异性参数化灵活性的强烈依赖。未发表的 2019 MMHT 分析~\cite{Thorne:2019mpt}
对 61 个 ATL7ZW 数据点得到 $\chi^2_E/N_\mathit{pt,E}\!=\! 1.76$，做法是：(a) 对单举带电流 (SI)DIS 截面纳入 NNLO 夸克质量修正
~\cite{Berger:2016inr,Gao:2017kkx}（轻微改善 DIS 与 ATL7ZW 数据集间的符合，参看 Sec.~\ref{sec:Qualitydimuon}）；(b) 对 $d,s$ 夸克使用灵活的 6 参数参数化。该 2019 MMHT 分析报告 $x\! >\! 10^{-3}$ 处 $s(x,Q)$ 较 MMHT'2014 显著增强，相应不确定度缩减。

NNPDF3.0 未主动拟合 ATL7ZW（其 $\chi_E^2/N_\mathit{pt,E}\!=\! 8.44$）；这些数据实现于 NNPDF3.1~\cite{Ball:2017nwa}，
得到 $\chi^2_E/N_\mathit{pt,E}\!=\! 2.14$ ——接近 CT18A/Z NNLO。
NNPDF3.1 观察到的奇异 PDF 增强与 CT18A NNLO 在 $x < 0.1$ 的结果相似。图~\ref{fig:CT18A2NNPDF31} 绘出 $Q\!=\!100$ GeV 处 CT18A 与 NNPDF3.1 NNLO 奇异 PDF 相对 CT18（不含 ATL7ZW 的基线）的比值。
在 CT18A 中，ATL7ZW 数据拉动最强的区域 $x\!\gtrsim\! 0.02$ 内我们看到高于 CT18 约 $\sim\!10\%$ 的超出量。尤其在这一 $x$ 区域，CT18A 中更大的 $s$-PDF 与 NNPDF3.1 拟合的奇异性紧密呼应。

\begin{figure}[t]
	\begin{tabular}{cc}
		\subfloat[$W^-$ production]{\includegraphics[width=0.35\textwidth]{images/chisq248WmCT18A_248_CT18A.pdf}}
		\subfloat[$W^+$ production]{\includegraphics[width=0.35\textwidth]{images/chisq248WpCT18A_248_CT18A.pdf}}
		\subfloat[$Z$ production]{\includegraphics[width=0.35\textwidth]{images/chisq248ZCT18A_248_CT18A.pdf}}
	\end{tabular}
	\caption{
		(a) 基于 CT18 与 CT18A/Z NNLO 的 ATL7ZW 数据理论预言比较。
		平移后数据（品红叉号）按 CT18A NNLO 计算。上面板给出微分截面随 $|\eta_l|$（带电流过程）或 $|y_{\ell \bar \ell}|$（中性流）的快度分布，下面板给出按 CT18A 理论归一的数据/理论比值。
	}
\label{fig:248_DT}
\end{figure}


\begin{figure}[p]
	\begin{tabular}{cc}
\subfloat[$W^-$ production]{\includegraphics[width=0.49\textwidth]{images/Residual_ATL7WZ_WmChart_shifted.pdf}}
\subfloat[$W^+$ production]{\includegraphics[width=0.49\textwidth]{images/Residual_ATL7WZ_WpChart_shifted.pdf}}\\
\subfloat[$Z$ production]{\includegraphics[width=0.49\textwidth]{images/Residual_ATL7WZ_ZChart_shifted.pdf}}
	\end{tabular}
	\caption{
		图~\ref{fig:248_DT} 所示三个 ATL7ZW 过程逐分箱的平移 $\mathrm{Theory}\!-\!\mathrm{Data}$ 残差，基于 CT18(A/Z) 计算。
	}
\label{fig:248_res}
\end{figure}

\begin{figure}[p]
\includegraphics[width=0.49\textwidth]{images/pResidual_CT18A_ATL7WZ.pdf}\quad \includegraphics[width=0.44\textwidth]{images/Nuisance_CT18A_ATL7WZ.pdf}\\
(a)\hspace{2in} (b)
	\caption{
		ATL7ZW 数据在 CT18A NNLO 下的 (a) 平移残差直方图与 (b) 讨厌参数直方图。
	}
\label{fig:248_hist}
\end{figure}

与 ATL7ZW 数据集的符合受 Drell-Yan 对产生 NNLO 代码选择的影响（讨论见 App.~\ref{sec:Appendix4xFitter} 末尾），也受 QCD 标度选择影响：CT18A/Z 拟合取 $\mu_{R,F}=Q$，MMHT（2019）取
$\mu_{R,F}=Q/2$，其中 $Q=M_W$（$M_{\ell\bar \ell}$）（$W$ 或 $Z$ 玻色子产生）。
在拟合中纳入伴随的 $W$ 玻色子-粲喷注（$W+c$）产生数据（如 NNPDF3.1 采用的 CMS 7 TeV 数据集~\cite{Chatrchyan:2013uja}）倾向于额外向上拉动 $s(x,Q)$。CT18A(Z) 尚未纳入 $W+c$ 测量，因为完整 NNLO 计算仍不可得。同样缺失的还有 SIDIS 双缪子产生微分截面的 NNLO 有质量重夸克贡献——重建被拟合 CCFR/NuTeV 截面时计算探测器接受度需要它。

\begin{figure}[tb]
\includegraphics[width=0.9\textwidth]{images/248rat_ifl5_lam_ct18Ann_lambda_q100.pdf}
	\caption{
		与上文 $L_2$ 敏感度图类似的图：显示 ATLAS $W/Z$ 数据关联系统不确定度在 CT18A NNLO 中对 $R_s(x,Q=100\,\mathrm{GeV})$ 的拉动，以相应讨厌参数平方的移动 $\Delta \lambda^2_\alpha$ 表示。
	}
\label{fig:248_l2nui}
\end{figure}

\subsubsection{与快度分布的比较}
\label{sec:ATL7ZWrap}

图s.~\ref{fig:248_DT}-\ref{fig:248_hist} 展示 CT18A/Z 预言相对 ATL7ZW 数据集的理论-数据比较、末态（赝）快度三个独立道的平移残差分布，以及最优拟合平移残差 $r_i$ 与讨厌参数 $\bar\lambda_\alpha$ 的累积直方图。

就单个数据点的描述而言，图形显示 CT18A/Z 预言能很好描述 $W^+$ 产生，但与 $W^-$ 产生在全轻子赝快度范围、以及与平均 $Z$ 玻色子快度为 0.9 和 1.3 分箱中的 $Z$ 产生都存在偏高的差异。

对这些数据的描述还需要五个讨厌参数 $\bar \lambda_\alpha$ 移动约 $\pm 2$ 个标准差——在 ATLAS
数据集中标记为 113、129、72、125 与 128。[这些参数混合了各种系统源的贡献，没有确定的物理解释。]

$L_2$ 敏感度技术的一个变体使我们得以证明：这些参数的变动与 ATL7ZW 测量所探 $x$ 区域内奇异性比 $R_s$ 的变化强烈相连。
前面的图中，$L_2$ 敏感度途径探讨的是单个实验的 PDF 变动与 $\Delta \chi^2_E$
的联系；但我们也可以计算 PDF 与单个讨厌参数贡献给 $\chi^2$ 的量之间的 $L_2$ 关联，即任意 $a\approx a_0$ 时按式~(\ref{Chi2a0l0}) 计入 $\Delta \chi^2_E(a)$ 的 $\Delta \bar \lambda^2_\alpha(a)$（$\alpha=1,...,\ N_\lambda$）。图~\ref{fig:248_l2nui} 以 ATL7ZW 数据为例演示这一点，展示 CT18A 中 $Q=100$ GeV 处 ATL7WZ 关联系统不确定度对 $R_s$ 的拉动。虽然大多数讨厌参数移动保持在 $|\Delta \lambda^2_\alpha|\! \sim\! 1\!-\!2$ 单位内，一小撮 $\bar \lambda_\alpha$
在不同 $x$ 处对被拟合 $R_s$ 高度敏感， notably 参数 113、72、120 与 129，其中一些几乎逼近（参数 72 与 129）或超过（参数 113 与 120）$|\Delta \lambda^2_\alpha|\! <\!2$ 的界限。


我们没有信息揭示驱动这些系统源的具体原因，但显然其中一些对图~\ref{fig:248_l2nui} 可读出的 $x$ 区间内偏好的 $R_s$ 行为具有深远影响。

\subsubsection{小结}

ATL7ZW 数据的精度使其在纳入 CT18 全球 PDF 拟合时对奇异夸克分布产生强烈影响。该数据在拟合中的存在（无论 CT18A 还是 CT18Z）使其 $\chi^2$ 相对直接用 CT18 PDF 计算大幅改善，但仍显著高于每自由度 1。该数据集的有效高斯变量也很大，只有 HERA I+II 可与之相比。总体而言，与其他 CT18 全球数据的张力导致拟合一致性下降（由强拟合优度判据~\cite{Kovarik:2019xvh} 定量刻画），并促成 ATL7ZW 数据进入 CT18A/Z PDF 而不进入 CT18 PDF。更多精密 LHC 数据或将进一步化解奇异性问题。在此之前，全面探索奇异性不确定度需要使用 CT18A 或 CT18Z PDF 组。


\subsection{CT18Z 中关于 PDF、$\alpha_s$ 与 $m_c$ 的扫描
\label{sec:LMCT18Z}
}

本附录以拉格朗日乘子扫描的若干结果作结。LM 扫描是在全球分析周期末端应用的强大技术，用以细察快速但近似的 Hessian 研究无法提供的精确概率分布。这里针对 CT18Z NNLO 给出的结果是对 Sec.~\ref{sec:QualityOverview} 就 CT18 所作同类讨论的补充。通过考察由完整拟合返回而非快速线性化近似给出的
$\Delta \chi^2$ 数值，我们能发现微妙的特征：实验间的张力，或某些 PDF 组合下拟合的不稳定/多解。

\begin{figure}[htbp]
\begin{center}
	\includegraphics[width=0.48\textwidth]{images/ATL7ZWWT_scan2Tct18z.pdf}\quad
	\includegraphics[width=0.48\textwidth]{images/NuTeVWT_scan2Tct18z.pdf}
	\caption{（左面板）在 CT18Z NNLO 拟合中变动 ATL7ZW 数据权重偏离默认拟合的 $w\!=\!1$ 时，主导数据集的总 $\chi^2$ 变化 $\Delta \chi^2$。（右面板）NuTeV $\nu,\, \bar{\nu}$ SIDIS 双缪子数据（Exp.~IDs=124、125）的类似图。
	}
\label{fig:lm_wht}
\end{center}
\end{figure}

\subsubsection{变动 ATLAS 7 TeV $W/Z$ 与 NuTeV 数据统计权重的扫描}
\label{sec:ATL7ZWchi2wt}

在 CT18Z NNLO 拟合内变动 ATL7ZW 数据权重~\cite{Collins:2001es}，考察若干最敏感实验 $\chi^2_E$ 的变化，即可审视 ATL7ZW 数据与其他实验间明显的张力。
结果见图~\ref{fig:lm_wht} 左面板：ATL7ZW 数据权重从 $w\!=\!0.1$（强烈淡化该信息，使 $\chi^2_E$ 移动 $\Delta \chi^2_E\! \sim\! +80$）
连续调到 $w\!=\!20$（强烈超重 ATLAS 数据，使 CT18Z 对其描述改善 $\Delta \chi^2_E\! \sim\! -40$）。表~\ref{tab:CMN} 报告的 ATL7ZW 数据 $\Delta \chi^2_E$ 不确定度对应总 $\chi^2$ 增加 36。

\begin{figure}[tb]
\begin{center}
	\includegraphics[width=0.48\textwidth]{images/Rs_scan2Tct18z_1.pdf}\quad
	\includegraphics[width=0.48\textwidth]{images/Rs_scan2Tct18z_2.pdf}\\
(a)\hspace{2.6in}(b)\\
	\caption{(a,b) CT18Z 拟合中 $Q=1.5$ GeV、$x=0.023$ 与 $x=0.1$ 处 $R_s$ 的拉格朗日乘子扫描，类比 CT18 的图~\ref{fig:LMRs}。
\label{fig:lm_rsz}}
\end{center}
\end{figure}

\begin{figure}[tb]
\begin{center}
	\includegraphics[width=0.6\textwidth]{images/Rs1_errorTct18a_T2MSTW.pdf}
	\caption{基于 $\chi^2$ 变动与 Hessian 估计得到的 $R_s(0.023,1.5\mbox{ GeV})$ 的 68\% C.L. 不确定度范围。对动态容忍度与个别实验的容忍度范围，绿色实（红色虚）带对应
\label{fig:rserrors}}
\end{center}
\end{figure}

图中还绘出当 ATL7ZW 数据集的 $\chi^2$ 权重改变时其 $\Delta \chi^2_E$ 变动最大的实验曲线。
左插图值得注意的特征是：反对以 5-10 权重纳入 ATL7ZW 数据集的，主要是 HERAI+II
单举 DIS 合并数据集（Exp.~ID=160），而非其他实验。除 LHCb 8 TeV $W/Z$
数据（250）——它在 ATL7ZW 权重增大时描述反而更好——外，其余所示实验都反对这种增大。这些实验包括 E866（203）与 NMC（104）$p/d$ 比、NuTeV 双缪子产生（124、125），以及 CMS 7 TeV $\mu$ 与 $e$ 不对称数据（266 与 267）。

当 ATLAS 数据被淡化（$w\!\sim\!0.1$）时，
NuTeV SIDIS 双缪子产生数据集（Exp.~IDs=124、125）的描述改善最多，进一步提示存在张力。

这些观察与右插图的伴随扫描一致：后者同样绘出对 NuTeV 双缪子数据权重变化最敏感实验的 $\Delta \chi^2$ 变化。特别考虑到所示实验的 $N_{pt}$，重度超重 NuTeV 数据会导致 ATL7ZW 数据点 $\chi^2$ 迅速恶化——实际上随 NuTeV 权重增大，其恶化快于完整 CT18Z 全球分析。事实上，当 NuTeV 权重增到约 5 时，单举 HERA（160）与 CCFR 中微子双缪子产生（126）的 $\chi^2_E$ 轻度改善；另一方面，E866 $pp$（204）、CCFR 反中微子双缪子（127）、CMS 8 TeV 电荷不对称（249）与 NMC $pd$
比（104）反对这种增大。


\subsubsection{CT18Z 奇异性比 $R_s$ 的 LM 扫描}
\label{sec:LMRsCT18Z}

图~\ref{fig:lm_rsz} 展示 CT18Z NNLO 在 $Q\!=\!1.5$ GeV、两个代表性动量分数
$x\!=\!0.023$（左面板）与 $x\!=\!0.1$（右）处的奇异性比 $R_s(x,Q)$ 拉格朗日乘子扫描，揭示出我们所做的其他 LM 扫描同样观察到的若干重要特征。

两面板均显示：HERA DIS 与 NuTeV/CCFR SIDIS 组等一方偏好较小的 $R_s$（在两个 $x$ 值处峰偏向 $R_s\!\sim\! 0.5$），而 ATL7ZW 数据正向拉动，并有较弱的拉动支持：E866 $pd$ 比（203）、BCDMS $F_2^d$（102），尤其 $F_2$ CCFR 中微子（110）数据——后者在 $x=0.1$ 非常显著。

当强迫 $R_s$ 取 0.7-0.8 以上数值时，两次扫描的 $\chi^2$
都开始不规则波动，且当 $R_s$ 被推至更高值时许多拟合不收敛。

鉴于 DIS 与 ATL7ZW 数据集拉动之间如此强烈的权衡，我们观察到：无论基于 CT 两级容忍度还是 MMHT 使用的动态容忍度，$R_s$ 的 Hessian 不确定度都不能完全捕获图~\ref{fig:lm_rsz} LM 扫描所揭示的真实 $\chi^2$ 行为。提醒读者：CTEQ-TEA Hessian 本征矢 PDF 采用两级方案构造~\cite{Lai:2010vv,Gao:2013xoa}，既防止全局 $\chi^2$ 增加过大，也防止单个实验 $\chi^2_E$ 增加过大。图~\ref{fig:rserrors} 示意 CT18A NNLO 拟合中 $R_s(0.023,1.5\mbox{ GeV})$ 68\% C.L. 不确定度区间的几种估计。[对 CT18Z NNLO 拟合发现类似。]最上面的估计基于 CT18A Hessian 本征矢组，其 68\% C.L. 估计由 90\% C.L. 非对称不确定度除以 1.645 得到。紧接 CT18A Hessian 估计之下的是左图~\ref{fig:lm_rsz} $R_s$ 扫描中 $\Delta \chi^2_{total}=10$ 对应的区间。可以看到：CT Hessian 不确定度名义上对应 $\Delta \chi^2_{total}\approx 36$，但其实际区间却与 LM 扫描真实 $\Delta \chi^2_{total}\approx 10$ 的区间相当。真实不确定度比 Hessian 估计暗示的更宽。

图~\ref{fig:rserrors} 下部，绿色实误差带给出单个实验 $\chi^2_E$ 分布的 68\% 分位数范围 $\xi_{68}$。红色虚误差带是这些 68\% 分位数范围经 $\chi^2_{E,0}/\xi_{50}$ 重标定的结果，$\chi^2_{E,0}$ 是最优全局拟合中实验 $E$ 达到的 $\chi^2$ 值。据此，第三行按 MSTW'2008 分析~\cite{Martin:2009iq} 第 6.2 节的程序构造 $R_s(0.023,1.5\mbox{ GeV})$ 68\% C.L. 不确定度的估计。所得区间可解读为 MSTW/MMHT 途径中沿 $R_s(0.023,1.5\mbox{ GeV})$ 梯度所选特定本征矢方向的动态容忍度。[其他本征矢与之垂直，对 $R_s(0.023,1.5\mbox{ GeV})$ 的不确定度无贡献。]

第三行的动态估计不确定度明显窄于其上两条不确定度区间。该动态不确定度等于 NuTeV 中微子与 ATL7ZW 数据集重标定 68\% 分位数之间的小重叠——前者从下方约束动态区间，后者从上方约束。若某些实验彼此强烈分歧，如此计算的 $R_s(0.023,1.5\mbox{ GeV})$ 动态不确定度将趋于过小。
同理，若某些实验强烈分歧，CT 分析的第二级惩罚在某些本征矢方向也可能导致过窄的不确定度。

\begin{figure}[tbp]
\centering
\includegraphics[width=0.49\textwidth]{images/alphas_scanTct18z_pjb05b.pdf}\quad
\includegraphics[width=0.46\textwidth]{images/alphas_scanTct18z_pjb05b2.pdf}\quad
\caption{CT18Z 强耦合常数在 $M_Z$ 标度的扫描，类比 CT18 的图~\ref{fig:lm_alphas}。同前，左面板给出对 $\alpha_s(M_Z)$ 具领先敏感度的若干实验的 $\Delta \chi^2$ 变化，右面板再次给出 CT18Z 所拟合全部实验的 $\chi^2$ 变化（分别归入 DIS、DY 与 top/jets 组合数据集）。同样地，右面板``Jets+top''曲线主要受喷注产生数据集影响。由于图~\ref{fig:lm_alphas} 所指出的 $t\bar{t}$ 数据的中等影响，导向 CT18Z 的 CT18 拟合变动使得 ATLAS 8 TeV $t\bar{t}$ 数据（Exp.~ID=580）此次入选左面板敏感实验组合。
\label{fig:lm_alphasz}}
\end{figure}


\subsubsection{CT18 关于 $\alpha_s(M_Z)$ 与 $m_c$ 的 LM 扫描}
\label{sec:CT18Zalphas}
CT18Z 的种种修改——尤其是在覆盖能量标度 $Q$ 最广的 DIS 实验中使用依赖
$x_B$ 的标度 $\mu_{F,x}$——改变了强耦合 $\alpha_s(M_Z)$ 的约束；扫描见图~\ref{fig:lm_alphasz}。

\begin{figure}[htbp]
\centering
\includegraphics[width=0.6\textwidth]{images/i2TZn3Q01-00mc1-35_DEchi2_re_ect.pdf}
\caption{
        类似图~\ref{fig:lm_mc}：用 CT18Z NNLO 按 Sec.~\ref{sec:AlphasDependence} 所述程序计算的粲极点质量 $m_c$ 数值扫描。
        }
\label{fig:lm_mcZ}
\end{figure}


尽管 $68\%$ C.L. 测定 $\alpha_s(M_Z)\! =\!
0.1169\! \pm\! 0.0027$ 与基于 CT18 的测定高度相符（后者仅稍弱、不确定度相同），实验拉动却有显著差异。最醒目的是：合并单举喷注产生数据与 Drell-Yan 实验的偏好出现分离——二者在 CT18 的图~\ref{fig:lm_mc} 中曾紧密对齐。合并 Drell-Yan 数据（含 ATL7ZW）如今与完整 CT18Z 拟合的偏好值非常接近，但 DIS 与喷注+$t\bar{t}$ 数据比 CT18 中更强地把 $\alpha_s(M_Z)$ 向相反方向拉动。

三类实验偏好的 $\alpha_s(M_Z)$ 区间对 DIS QCD 标度如此明显的依赖，或表明存在未被名义 68\% C.L. 不确定度 0.0027 计及的重要 NNLO 之外的不确定度。

类似地，替代性 CT18X/A/Z 全球分析中的选择可使这些拟合对粲极点质量 $m_c^{pole}$ 产生不同的偏好。
虽然 Sec.~\ref{sec:AlphasDependence} 已就 CT18 详细描述了 $m_c$ 扫描，对 CT18Z NNLO 重复该扫描得到的行为略有不同，见图~\ref{fig:lm_mcZ}。尽管 CT18Z 最终得到十分相近的 $m_c$ 中心值，敏感实验之间的相互作用与图~\ref{fig:lm_mc} 所示 CT18 的情形略有不同。合并 HERA 数据（160）对 $m_c$ 的拉动在 CT18Z 减小，HERA 粲产生数据（Exp.~ID=147）的偏好 $m_c$ 则略升。此外，ATL7ZW 数据也对较大粲质量表现出适度偏好；后两者共同使 CT18Z 相对 CT18 的 $m_c$ 中心偏好值小幅升高，而不确定度范围相近。



\clearpage

\section{CT18 拟合优度函数与关联不确定度的处理}
\label{sec:chi2_app}
%
%
\newcommand{\betamat}{b}

本附录总结 CT18 族拟合中拟合优度函数 $\chi^2$ 的实现与讨厌参数的边缘化。对后者，通常沿用自 CTEQ6 分析起采纳的程序：利用实验发表的关联矩阵估计关联不确定度。对少数不提供关联矩阵的数据集，我们把协方差矩阵写成区分非关联与关联成分的近似形式较为便利。转换算法在本附录末说明。


{\bf 定义。}
Sec.~\ref{sec:chi2} 式~(\ref{Chi2sys}) 引入了近期 CT 拟合使用的标准拟合优度函数 $\chi^{2}$。此处用矩阵记号回顾其处理。

我们对每个 $k$ 用 $s_{k}$ 为单位表达 $D_{k},$ $T_{k}$ 与 $\beta_{k\alpha}$，即引入长度 $N_{pt}$ 的向量
$d\equiv S^{-1}\left(D-T(a)\right)$，其中
\begin{equation}
S\equiv\mbox{diag}\left\{ s_{1},s_{2},...,s_{N_{pt}}\right\} .\label{S}
\end{equation}
类似地，$\lambda\equiv\{\lambda_{\alpha}\}$ 是长度
$N_{\lambda}$ 的向量；$\betamat\equiv S^{-1}\beta$ 是维数 $N_{pt}\times N_{\lambda}$ 的矩形矩阵。在该矩阵记号下，式~(\ref{Chi2sys}) 写成
\begin{equation}
\chi_{E}^{2}(a,\lambda)=\left(d-\betamat\lambda\right)^{T}(d-
\betamat\lambda)+\lambda^{T}\lambda.\label{Chi2sys2}
\end{equation}

{\bf $\chi^{2}$ 的极小化。}
$\chi^{2}$ 极小值（相对讨厌参数）的解借助以下矩阵求出：
\begin{align}
\mathcal{A}\equiv\mathcal{I}+ {\betamat}^{T}\betamat;~~~ & C\equiv I+\betamat\betamat^{T};\label{ACm1}\\
\mathcal{A}^{-1}=\mathcal{I}-\betamat^{T}C^{-1}\betamat;~~~ & C^{-1}=I-\betamat\mathcal{A}^{-1}\betamat^{T}.\label{Am1C}
\end{align}
大写罗马字母与花体字母分别表示维数 $N_{pt}\times N_{pt}$
与 $N_{\lambda}\times N_{\lambda}$ 的矩阵。因此 $I$
是 $N_{pt}\times N_{pt}$ 单位矩阵，$\mathcal{I}$ 是
$N_{\lambda}\times N_{\lambda}$ 单位矩阵。$C$ 与 $\mathcal{A}$
分别是数据值空间与讨厌参数空间中的协方差矩阵（适当以 $s_{k}$ 归一）。式~(\ref{Am1C}) 中 ${\cal A}^{-1}$ 与 $C^{-1}$ 的关系可用 $\betamat^{T}\betamat=\mathcal{A}-\mathcal{I}$ 与
$\betamat\betamat^{T}=C-I$ 证明。协方差矩阵对称：$\mathcal{A}^{T}=\mathcal{A},$
$C^{T}=C$。

对每个理论输入 $a$，$\chi^{2}$ 可解析地对 $\lambda$ 极小化，即求解 $\partial\chi^{2}/\partial\lambda_{\alpha}=0$
~\cite{Pumplin:2002vw}。解为
\begin{align}
 & \bar{\lambda}(a)=\mathcal{A}^{-1}\betamat^{T}d=\betamat^{T}C^{-1}d;\label{lambda0}\\
 & \bar{\lambda}^{T}(a)=d^{T}\betamat\mathcal{A}^{-1}=d^{T}C^{-1}\betamat.
\end{align}
 所有实验 $E$ 的全局极小 $\chi^{2}(a_{0},\lambda_{0})$ 可数值求得：
\begin{equation}
\chi^{2}(a_{0},\lambda_{0})=\sum_{E}\ d_{0}^{T}C^{-1}d_{0},\label{Chi2CovMat}
\end{equation}
 其中 $d_{0}\equiv S^{-1}\left(D-T(a_{0})\right)$，$\lambda_{0}\equiv\bar{\lambda}(a_{0})$。

该方程还有一个等价形式：
\begin{equation}
\chi^{2}(a_{0},\lambda_{0})=\sum_{E}\ \left(r_{0}^{T}r_{0}+\lambda_{0}^{T}\lambda_{0}\right),\label{Chi20rl}
\end{equation}
其中 $r_{0}$ 是最优拟合平移残差：
\begin{equation}
r_{0}\equiv S^{-1}\left(d_{0}-\betamat\lambda_{0}\right)=\bar{r}(d_{0}),\mbox{~with~}\bar{r}(d)\equiv C^{-1}d.\label{r0}
\end{equation}
表示 (\ref{Chi20rl}) 特别富有启发性。我们预期，在理论与实验 $E$ 的良好拟合中，量化与单个数据点符合程度的平移残差 $r_{0}$，以及量化系统位移的讨厌参数 $\bar{\lambda}_{0}$，应服从各自的标准正态分布 ${\cal N}(0,1)$。把两个经验分布与预期的 ${\cal N}(0,1)$ 分布比较，即可作为拟合优度与系统误差实现的检验~\cite{Kovarik:2019xvh}。

{\bf 协方差矩阵的分解。}
在我们分别已知非关联与完全关联系统误差这一常见情形中，式~(\ref{Chi2CovMat1}) 的 $\chi^2$ 形式与对给定 $a$ 优化讨厌参数 $\lambda$ 所得式~(\ref{Chi2CovMat}) 一致。此时认定
\begin{equation}
\mbox{cov}=SCS,~~~\mbox{cov}^{-1}=S^{-1}C^{-1}S^{-1}.\label{SCS}
\end{equation}
 矩阵元按式s.~(\ref{ACm1}) 与 (\ref{Am1C}) 为
\begin{align}
\left(\mbox{cov}\right)_{ij} & =s_{i}^{2}\delta_{ij}+\sum_{\alpha=1}^{N_{\lambda}}\beta_{i\alpha}\beta_{j\alpha},\label{cov}\\
\left(\mbox{\mbox{cov}}^{-1}\right)_{ij} & =\frac{\delta_{ij}}{s_{i}^{2}}-\sum_{\alpha_1,\alpha_2=1}^{N_{\lambda}}\frac{\beta_{i\alpha_1}}{s_{i}^{2}}{\cal {\cal A}}_{\alpha_1\alpha_2}^{-1}\frac{\beta_{j\alpha_2}}{s_{j}^{2}}.\label{covm1}
\end{align}
特别地，对角元 $(\mbox{cov})_{ii}$［不求和］是统计、非关联系统与关联系统不确定度的平方和：
\begin{equation}
(\mbox{cov})_{ii}=s_{i,\mbox{\scriptsize stat}}^{2}+s_{i,\mbox{\scriptsize uncor sys}}^{2}+s_{i,\mbox{\scriptsize cor sys}}^{2},\label{s2i}
\end{equation}
其中 $s_{i,\mbox{\scriptsize cor sys}}^{2}\equiv\sum_{\alpha}\beta_{i\alpha}^{2}$。
借助式~(\ref{SCS})，式~(\ref{r0}) 的平移残差写作
\begin{equation}
r_{i}\equiv\frac{D_{i}^{sh}(a)-T_{i}(a)}{s_{i}}=s_{i}\sum_{j=1}^{N_{pt}}\left(\mbox{\mbox{cov}}^{-1}\right)_{ij}\left(D_{i}-T_{i}(a)\right).\label{ricov}
\end{equation}
因此计算它们需要知道完整的非关联误差 $s_{i}=\sqrt{s_{i,\mbox{\scriptsize stat}}^{2}+s_{i,\mbox{\scriptsize uncor sys}}^{2}}$。

{\bf 由协方差矩阵求关联矩阵。}某些实验测量（如 LHCb 8 TeV W/Z 产生~\cite{Aaij:2015zlq}）只提供完整协方差矩阵，使按式~(\ref{ricov}) 直接重构平移残差不可能。在若干情形中，我们发现可行的做法是：假定系统移动由某数目 $M_{\lambda}$（$M_{\lambda}\leq N_{\lambda}$）个完全关联线性组合主导，从而由所给协方差矩阵 $\mbox{cov}\equiv K_{0}$ 迭代重构近似的非关联系统与关联系统贡献 $s_{k,\mbox{\scriptsize uncor sys}}$ 与 $\beta_{k\alpha}$。该近似利用协方差矩阵及其对角元的正定性，参看式~(\ref{s2i})。它用一个数值上相近的矩阵代表原协方差矩阵：后者是一个对角矩阵 $\Sigma$（解读为总非关联误差之和）与一个非对角方阵 $K$（解读为关联矩阵 $\beta$ 及其转置之积）之和。

具体地，
设求得 $K_{0}$ 的本征值 $x_{k}^{2}$ 并按降序排列：
\begin{align}
K_{0} & =O^{T}xO,\mbox{ with }x=\mbox{diag}\left\{ x_{1}^{2},x_{2}^{2},...,x_{N_{pt}}^{2}\right\} ,\label{K0}\\
 & x_{1}^{2}\geq x_{2}^{2}\geq...\geq x_{N_{pt}}^{2}>0.
\end{align}
其中 $O$ 是正交矩阵。把 $x$ 划分为含最大 $M_{\lambda}$
个本征值的矩阵 $y$ 与含最小 $(N_{pt}-M_{\lambda})$
个者的矩阵 $z$：
\begin{align}
y & \equiv\mbox{diag}\left\{ x_{1}^{2},x_{2}^{2},...,x_{M_{\lambda}}^{2},0,...,0\right\} ,\\
z & \equiv\mbox{diag}\left\{ 0,...,x_{M_{\lambda+1}}^{2},...,x_{N_{pt}}^{2}\right\} .\label{z}
\end{align}
注意到矩阵 $Y=O^{T}yO$ 与
$Z=O^{T}zO$ 的对角元非负，我们随后构造对角矩阵
$\Sigma_{1}\equiv\mbox{diag}\left\{ Z_{11},...,Z_{N_{pt}N_{pt}}\right\} $
（$Z_{ii}>0$），以及另一个正定矩阵 $K_{1}\equiv Y+Z-\Sigma_{1}$。
可在第 $a$ 步计算 $\Sigma_{a+1}$ 与 $K_{a+1}$ 来迭代式s.~(\ref{K0})-(\ref{z}) 的步骤：
\begin{align}
\Sigma_{a+1} & \equiv\Sigma_{a}+\mbox{diag}\left\{ Z_{11},...,Z_{N_{pt}N_{pt}}\right\} ,\\
K_{a+1} & \equiv Y+Z-\mbox{diag}\left\{ Z_{11},...,Z_{N_{pt}N_{pt}}\right\} .
\end{align}
其中矩阵 $Y,Z$ 每步都以 $K_{a}$ 为输入重新计算。经过足够多步 $a$ 之后，和
\begin{equation}
C_{a}\equiv\Sigma_{a}+Y_{a}\label{Ca}
\end{equation}
趋近一个在 $L_{p}$ 范数 $\sum_{i,j=1}^{N_{pt}}\left|(K_{0})_{ij}-(C_{a})_{ij}\right|^{p}$（$p=2$ 或 $1$）意义上接近输入矩阵
$K_{0}$ 的渐近矩阵。若非关联与完全关联成分的提取可行，则选取足够大的 $M_{\lambda}$ 即可使渐近 $L_{p}$ 距离很小。对 CT18 数据集中仅提供协方差矩阵的三个实验~\cite{Aaij:2015gna,Aaij:2015zlq,Khachatryan:2016pev}，收敛良好的 $M_\lambda$ 值介于 $N_{\lambda}/2$ 与 $N_{\lambda}$ 之间。

比较式s.~(\ref{cov}) 与 (\ref{Ca})，我们认定
\begin{align}
(\Sigma_{a})_{ij} & \approx s_{i}^{2}\delta_{ij},\\
(Y_{a})_{ij} & \approx\sum_{\alpha=1}^{M_{\lambda}}\beta_{i\alpha}\beta_{j\alpha}.
\end{align}
因此 $s_i$ 由 $(\Sigma_a)_{ij}$ 估计；$\beta_{i\alpha}$ 可由 $(Y_{a})_{ij}$ 经奇异值分解估计。


\section{非微扰参数化形式}
\label{sec:AppendixParam}

如 Sec.~\ref{sec:Paramstudies} 所述，为现实地估计参数化 PDF 不确定度，CT 全球分析在起始标度 $Q\! =\! Q_0$ 探索宽范围的 PDF 参数化形式。
这些研究的目标是 ({\bf i}) 挑选足够灵活、能在不发生过拟合的前提下拟合庞大高能数据集的函数形式；({\bf ii}) 理解与参数化选择相联系的不确定度。

CT14 论文~\cite{Dulat:2015mca} 的附录阐述了指导参数化形式选择（包括 CT18 分析采用的那些）的主要理据。初始标度 $Q_0$ 处以自由参数 $a_k$ 表示的一般函数形式为
\begin{equation}
    f_i(x,Q_0) = a_0 x^{a_1-1} (1-x)^{a_2} P_i(y; a_3, a_4,...).
    \label{eq:fiQ0}
\end{equation}
系数 $a_1$、$a_2$ 控制 $f_i(x,Q_0)$ 在极限 $x\rightarrow 0$ 与 $1$ 的渐近行为。$P_i(y;a_3,a_4, ...)$ 是依赖于 $y=f(x)$（如 $y\equiv \sqrt{x}$）的 Bernstein 多项式之和（亦称 B\'ezier 曲线），在整个区间 $0<x<1$ 上极为灵活。中间阶段尝试过多种函数形式 $P_i(x; a_3, a_4,...)$；在名义参数化中，价 PDF 共用同一组参数 $a_1$，海 PDF 另用一组。该选择保证相应 PDF 之比在极限 $x\to 0$ 趋于有限值。我们也对 $x\to 1$ 的参数 $a_2$ 施加类似关系。最后，为减少 $(1-x)^{a_2}$ 中系数 $a_2$ 与 $P_i(x;a_3,...)$ 其余系数在 $x\rightarrow 1$ 处的伪关联，我们把 $P_i(x;a_3,...)$ 中某些参数用其他参数表达，消去其中的线性 $(1-x)$ 项，即要求
\begin{equation}
    \lim_{x\to 1} xf_i(x,Q_0) = a_0 (1-x)^{a_2} \left(1 + {\cal O}\left((1-x)^2\right)\right).
    \label{PixTo1}
\end{equation}
在 CT18 中，该手续引出式s.~(\ref{eq:a6a1}) 与 (\ref{eq:a5a3}) 中的关系。对价夸克，这使我们能用四个自由参数（而非五个或更多）取得良好的 $\chi^2$ 值。完整手续围绕文献~\cite{Dulat:2015mca} 式 (A.16) 解释。

\begin{table}[tb]
\begin{tabular*}{\textwidth}{c| @{\extracolsep{\fill}} cccccc}
\hline
best-fit parameters, &&&&&&\\
  {\bf CT18} &  $u_v$          &   $d_v$          &    $g$            &   $u_\mathrm{sea}=\bar{u}$       &      $d_\mathrm{sea}=\bar{d}$       &    $s=\bar{s}$         \tabularnewline
\hline
\% mom. fraction     & 32.5    & 13.4   & 38.5 & $2.8$ & $3.6$ & $1.3$ \tabularnewline
$a_0$                & $3.385^\mathrm{SR}$    & $0.490^\mathrm{SR}$   & $2.690$ & 0.414 & 0.414 & $0.288$ \tabularnewline
$a_1$                                  &  $0.763$     &    $0.763$      &    $0.531$     &    $-0.022$      &      $-0.022$      &    $-0.022$     \tabularnewline
$a_2$                                  &  $3.036$     &    $3.036$      &    $3.148$     &    $7.737 $      &      $7.737$       &    $10.31$     \tabularnewline
$a_3$                                  &  $1.502$     &    $2.615$      &    $3.032$     &    $(4)$          &      $(4)$          &    $(4)$         \tabularnewline
$a_4$                                  &  $-0.147$    &    $1.828$      &    $-1.705$    &    $0.618 $      &      $0.292$       &    $0.466$      \tabularnewline
$a_5$                                  &  $1.671$     &    $2.721$      &      ---        &    $0.195 $      &      $0.647$       &    $0.466$      \tabularnewline
$a_6$                                  &     ---       &       ---        &      ---        &    $0.871 $      &      $0.474$       &    $0.225$      \tabularnewline
$a_7$                                  &     ---       &       ---        &      ---        &    $0.267 $      &      $0.741$       &    $0.225$      \tabularnewline
$a_8$                                  &     ---       &       ---        &      ---        &    $0.733 $      &        (1)          &      (1)         \tabularnewline
\hline
\hline   \\
{\bf CT18Z}          &  $u_v$          &   $d_v$          &    $g$            &   $u_\mathrm{sea}=\bar{u}$       &      $d_\mathrm{sea}=\bar{d}$       &    $s=\bar{s}$            \tabularnewline
\hline                                                                              \% mom. fraction     & 31.8    & 13.1   & 38.2 & $2.9$ & $3.7$ & $1.9$ \tabularnewline
$a_0$                & $3.631^\mathrm{SR}$    & $0.254^\mathrm{SR}$   & $0.668$ & 0.519 & 0.519 & 0.524 \tabularnewline

$a_1$                                  &  $0.787 $    &    $0.787$      &    $0.289 $    &    $0.0096 $      &      $0.0096$       &    $0.0096 $     \tabularnewline
$a_2$                                  &  $3.148 $    &    $3.148$      &    $1.872 $    &    $8.27 $      &      $8.27$       &    $11.4$     \tabularnewline
$a_3$                                  &  $1.559 $    &    $3.502$      &    $3.538 $    &    $(4)$          &      $(4)$          &    $(4)$         \tabularnewline
$a_4$                                  &  $-0.075$    &    $1.865$      &    $-1.665$    &    $0.679 $      &      $0.300$       &    $0.653 $     \tabularnewline
$a_5$                                  &  $1.605 $    &    $3.599$      &      ---        &    $-0.0016$      &      $0.532$       &    $0.653 $     \tabularnewline
$a_6$                                  &     ---       &       ---        &      ---        &    $1.085 $      &      $0.753$       &    $0.054 $     \tabularnewline
$a_7$                                  &     ---       &       ---        &      ---        &    $0.045 $      &      $0.440$       &    $0.054 $     \tabularnewline
$a_8$                                  &     ---       &       ---        &      ---        &    $0.759 $      &        (1)          &      (1)         \tabularnewline     %
\hline
\end{tabular*}\caption{
	$Q_0=1.3 $ GeV 处动量分数百分比与 CT18、CT18Z NNLO 拟合 PDF 的最优拟合参数值。动量分数百分比按各部分子味所列参数评估为 $\langle x \rangle_{f}$；进入动量
	求和规则的海夸克总动量百分比应为上列 $u_\mathrm{sea}$、$d_\mathrm{sea}$、$s=\bar{s}$ 数值的两倍。各参数化对应的函数形式显式定义于本节正文式s.~(\ref{eq:v_param})--(\ref{eq:ubar_exp})。某 PDF 味未被主动拟合的参数条目以短横``---''标注；括号内数值为固定参数。
	如上方 $\#^\mathrm{SR}$ 注记所示，$u_v,d_v$ 的归一化 $a_0$ 由价求和规则（SR）导出；其余归一化则由作为自由参数拟合的 $\langle x\rangle_g$ 与 $\langle x\rangle_{\bar s+\bar s}/\langle x\rangle_{\bar u+\bar d}$，加上用动量求和规则算得的
	$\langle x\rangle_{\bar u+\bar d+\bar s}$ 导出。
}
\label{tab:parameters}
\end{table}

在 CT18 中，名义非微扰参数化大体类似于作为 \CTHERAII{} 基础的那一套，但对海夸克分布采用了略更灵活的参数化。这种增强的灵活性因纳入对核子海夸克含量具有直接敏感度的 LHC Run-1 数据而被 necessitated。
参数化起始标度价分布 $u_v$ 与 $d_v$ 的函数形式是参数
$y\! \equiv\! \sqrt{x}$ 的多项式：
\begin{eqnarray}
\label{eq:v_param}
q_v(x,Q=Q_0) &=& a_0\, x^{a_1-1}(1-x)^{a_2}\,P^v_a(y), \\
P^v_a(y) &=& \sinh{[a_3]}         (1-y)^4 +
\sinh{[a_4]}   4 y   (1-y)^3 +
\sinh{[a_5]}   6 y^2 (1-y)^2 +
a_6   4 y^3 (1-y)   +
y^4,             \nonumber
\end{eqnarray}
其中
\begin{eqnarray}
a_6 &=& 1 + \frac{1}{2} a_1\  \label{eq:a6a1}
\end{eqnarray}
以满足式~(\ref{PixTo1})。

我们强调：尽管 $u_v$ 与 $d_v$ 分布的非微扰形式相同，这两味的 $P^v_a(y)$ 参数是分开拟合的；不过我们施加约束 $a^{u_v}_1\! =\! a^{d_v}_1$ 与
$a^{u_v}_2\! =\! a^{d_v}_2$ 于前置因子指数，以保证味比行为良好并与 Regge 期望及夸克计数规则在极限 $x\! \to\! 0,1$ 一致。

对胶子 PDF，我们同样拟合 $y\! =\! \sqrt{x}$ 的多项式，但形式为
\begin{eqnarray}
g(x,Q=Q_0) &=& a_0\, x^{a_1-1}(1-x)^{a_2}\,P^g_a(y), \\
P^g_a(y) &=& \sinh{[a_3]}         (1-y)^3 +
\sinh{[a_4]}   3 y   (1-y)^2 +
a_5   3 y^2 (1-y)   +
y^3, \nonumber
\end{eqnarray}
其中参数 $a_5$ 固定为
\begin{eqnarray}
a_5 &=& ( 3 + 2 a_1 ) \big/ 3. \label{eq:a5a3}
\end{eqnarray}

对轻夸克海的分布，我们相对 CT14 采用稍更灵活的形式。在这些情形中，对 $\bar{u}$、$\bar{d}$ 与 $\bar{s}=s$ 使用 $y\! \equiv\! 1\! -\! (1\! -\! \sqrt{x})^{a_3}$ 的多项式，三条海分布统一固定 $a_3\! =\! 4$。
海夸克 PDF $\bar{q}(x,Q_0)$ 参数化为
\begin{eqnarray}
\label{eq:s_param}
\bar{q}(x,Q=Q_0) &=& a_0\, x^{a_1-1}(1-x)^{a_2}\,P^{\bar{q}}_a(y), \\
P^{\bar{q}}_a(y) &=&             (1-y)^5 +
a_4   5 y   (1-y)^4 +
a_5  10 y^2 (1-y)^3 +
a_6  10 y^3 (1-y)^2   \nonumber \\
&+& a_7   5 y^4 (1-y)   +
a_8     y^5\ . \nonumber
\end{eqnarray}
完整的参数形式（式~(\ref{eq:s_param})，固定 $a_3 = 4$）用于 $\bar{u}$-PDF；对 $\bar{d}$ 固定 $a_8=1$。由于对核子奇异性的经验约束相对缺乏，CT18(Z) 中 $s(x,Q)$ 的上述参数并非全部自由浮动：我们对奇异性参数施加约束 $a_4\! =\! a_5$、$a_6\! =\! a_7$ 且 $a_8=1$。

与价分布一样，我们约束前置因子指数
\begin{equation}
	a^{\bar{u}}_1 = a^{\bar{d}}_1 = a^s_1\ ,
\label{eq:ubar_exp}
\end{equation}
以保证奇异性压低比 $R_s\! =\! [s+\bar{s}] \big/ [\bar{u}+\bar{d}]$ 在极限 $x\!\to\!0$ 取有限值，以及 $\int^1_0 dx\, [\bar{d}-\bar{u}]$ 收敛。我们绑定 $\bar{u}$-与 $\bar{d}$-分布的高 $x$ 指数 $a^{\bar{u}}_2\! =\! a^{\bar{d}}_2$，以稳定 $x\!\to\!1$ 的 $\bar{d}/\bar{u}$。各海夸克 PDF 的归一化由价夸克与动量求和规则、以及作为自由参数拟合的一阶矩 $\langle x \rangle_{g}$ 与比值 $\langle x\rangle_{\bar s+\bar s}/\langle x\rangle_{\bar u+\bar d}$ 计算。由于参数化不能决定奇异-非奇异 PDF 之比，我们施加适当的拉格朗日乘子约束，把 $\left(s(x,Q_0) + \bar s(x,Q_0)\right)/\left(\bar u(x,Q_0) + \bar d(x,Q_0)\right)$ 限制在 $x=10^{-8}$ 处的区间 $[0.2, 2.0]$ 与 $x=10^{-5}$ 处的区间 $[0.4, 1.8]$。

表~\ref{tab:parameters} 汇总上述参数对 CT18（上行）与 CT18Z（下行）的中心拟合值。标记``---''的参数条目不作为相应味的自由度参与拟合。

\section{拟合代码的发展}
\label{sec:AppendixCodeDevelopment}
把超过 10 个新 LHC 数据集纳入 CT18 分析，需要对 CTEQ 全球分析软件进行大量升级。
硬散射截面计算使用各种快速接口成为必需与惯例。
即便实现了 \texttt{ApplGrid} 等快速接口之后，精密全球拟合仍是耗时任务，原因在于实验数据规模巨大，且需探索多参数概率分布来估计 PDF 不确定度。从 CT14 过渡到 CT18 时，CTEQ-TEA 代码经过修订，把过去串行执行的多种操作并行化。

典型 CTEQ-TEA 全球分析的计算过程可分为三个阶段（``层''），见图~\ref{fig:gfit}。

外层 LY0 对应以不同输入或约束重复全球拟合，例如不同的 QCD 耦合、重夸克质量、非微扰函数形式或拉格朗日乘子约束。

中层 LY1 对应单次这样的拟合：程序扫描 PDF 参数空间并为固定的输入组合构建概率（$\chi^2$）分布。
该步骤提供量化参数空间不确定度的最优拟合与误差 PDF 组。

\begin{figure}[tb]
	\begin{center}
		\includegraphics[width=0.99\textwidth]{images/CT18code.pdf}
	\end{center}
	\vspace{-2ex}
	\caption{\label{fig:gfit}
		CT18 全球分析流程图，由外层（LY0）、
		中层（LY1）与核心过程（LY2）组成。
	}
\end{figure}

LY1 层内，全局 $\chi^2$ 的计算（即``层 LY2''）是拟合的核心部分，将对每一组 PDF 参数组合重复执行。
LY2 为数千个数据点计算截面并对每个数据集计算各自的 $\chi^2$。


图~\ref{fig:gfit} 所示的大部分计算工作都可并行化。
例如在 LY0 阶段，可以直接向大型计算集群提交数千个输入各异的并行拟合，因为这些拟合相互独立。为获得 CT18(Z) 分析所呈现的多样化结果系列（包括计算昂贵的拉格朗日乘子扫描），使用了 MSU 与 SMU 的高性能计算集群。

在 LY1，主要任务是在大维度参数空间（$N_{par}\sim 30$）寻找 $\chi^2$ 全局极小。
合适并行技术的选择强烈依赖于极小化算法。
CT18 分析采用 \texttt{MINUIT} 软件包~\cite{James:2004xla} 实现的变尺度梯度下降法。
它涉及 $\chi^2$ 一阶与
二阶导数的数值计算，并结合沿 PDF 参数空间固定方向的顺序极小搜索。
导数计算高度可并行化：CPU 时间开销分别大致按 $1/N_\mathit{par}$ 与 $2/(N^2_\mathit{par}\!+\!N_\mathit{par})$ 随参数个数 $N_{par}$ 标度。

在核心部分，对不同数据集（含其截面）各自 $\chi^2$ 的计算如今也是同时进行的。此阶段 \texttt{MPI} 或 \texttt{OpenMP} 并行协议皆适用；后者局限于共享内存平台，但对我们的拟合代码所需改动较少。
具体地，我们采用类似 \texttt{OpenMP} 的途径以缩小拟合代码内部的改动范围。
计算实验的 $\chi^2_E$ 值时，Linux 的 \verb|fork| 命令把主程序分裂为多个线程，
每个线程持有主存副本并独立执行计算。
随后 \verb|join| 命令收集各线程的结果并送回主程序。
这一 \verb|fork-join| 算法的实现借鉴了广泛用于多线程蒙特卡洛积分的
CUBA 库~\cite{Hahn:2004fe}。

\section{ATLAS 喷注产生截面数据的去关联}
\label{sec:ATLASjetdecorrel}
在我们的全球分析以及其他分析~\cite{Harland-Lang:2017ytb} 中均已观察到：要对 Run-1 LHC 喷注产生数据建立稳健的理论描述，一般需要数值手段对部分关联系统不确定度去关联。我们的工作遵循实验合作组关于此类去关联程序的建议。

例如，ATLAS 7 TeV 单举喷注产生数据~\cite{Aad:2014vwa}（Exp.~ID=544）开箱即用时 $\chi^2$ 高得不可接受，而应用 ATLAS 合作组提出、汇总于文献~\cite{Aaboud:2017dvo} 附录的部分去关联选项后，数据与 NNLO 理论的符合得到改善。
CT18(Z) 拟合实现了以 $R=0.6$ anti-$k_t$ 喷注算法定义的 ATLAS 单举喷注截面。对该数据集，我们检验了文献~\cite{Aaboud:2017dvo} 表 6 总结的去关联程序，
发现其中一些对 ATLAS 7 TeV 数据的 $\chi^2_E/N_\mathit{pt,E}$ 有实质影响。最终我们遵循 ATLAS 实验家自己的具体建议~\cite{Bogdan}：对两项喷注能量标度（JES）不确定度去关联，即 JES MJB fragmentation（JES16）与 JES flavor response
（JES62）。按文献~\cite{Aaboud:2017dvo} 表 6 上半部分、分别依表 4 的 Option 17 与 14 对 JES16 与 JES62 去关联，我们获得了最大的 $\chi^2$ 缩减。

文献~\cite{Aaboud:2017dvo}
还详述了实验模拟中选定不确定度的去关联选项；
其中我们专门考察了对非微扰修正（`eNPC'）相关联误差去关联的效应，但按 $\chi^2$ 衡量其对 \CTHERAII{} 拟合质量的影响可忽略。

我们实现的 JES 去关联的做法是：把给定关联不确定度拆分为 2-3 个子误差，使其平方和复原原关联不确定度。该手续依赖分箱；所得去关联误差取决于喷注快度与横动量。例如我们把 JES MJB fragmentation 误差 $\delta_{16}$ 去关联为
三个分量：
\begin{align}
	\delta^a_{16} &= \delta_{16}\, \Big\{ 1 - L^2\big( \ln(p_T[\mathrm{TeV}]),\ln(0.1),\ln(2.5) \big) \Big\}^{1/2}\, \left( 1 - L^2\big(|y|,0,1\big) \right)^{1/2}, \nonumber \\
	\delta^b_{16} &= \delta_{16}\, \Big\{ 1 - L^2\big( \ln(p_T[\mathrm{TeV}]),\ln(0.1),\ln(2.5) \big) \Big\}^{1/2}\, L\big(|y|,1,3\big), \nonumber \\
	\delta^c_{16} &= \big( \delta^2_{16} - (\delta^a_{16})^2 - (\delta^b_{16})^2 \big)^{1/2}\ ,
\label{eq:dec16}
\end{align}
其中
\begin{equation}
L(x,x_\mathit{min},x_\mathit{max}) \equiv \frac{x - x_\mathit{min}}{x_\mathit{max} - x_\mathit{min}}\ \iff x \in [x_\mathit{min},x_\mathit{max}]\ ,
\end{equation}
否则当 $x < x_\mathit{min}$ 时 $L = 0$，$x > x_\mathit{max}$ 时 $L = 1$。注意式~(\ref{eq:dec16}) 支配去关联不确定度的大小，而其符号继承原不确定度。类似算法应用于 JES flavor response JES62 的去关联。

类似的考虑也应用于 CMS 7 与 8 TeV 单举喷注产生数据集（实验 542 与 545）。具体而言，对 7 TeV CMS 喷注数据，应用文献~\cite{Khachatryan:2014waa} 的去关联方法处理 JEC2，并对 $2.5 < |y| < 3.0$ 分箱附加一项去关联\footnote{私人通信，M.~Voutilainen。}以得到 6 个子不确定度。对 8 TeV CMS 喷注数据，系统不确定度按文献~\cite{Khachatryan:2016kdb} 的处理用 xFitter 框架获得。
在此基础上，我们发现来自 \texttt{NNLOJet}~\cite{Ridder:2015dxa,Gehrmann-DeRidder:2017mvr,Currie:2016bfm,Currie:2017ctp} 现有 NNLO 预言蒙特卡洛积分的残余涨落会抬高单举喷注产生实验的 $\chi^2$，等价地说，若用光滑函数拟合制表的 NNLO/NLO 修正，会带来某种不确定度。
在 CT18 分析中，对三个喷注实验的蒙特卡洛（MC）理论不确定度通过添加 0.5\% 的整体{\it 非关联}不确定度（``MC unc.''）估计——这是与 NNLO/NLO $K$ 因子的蒙特卡洛生成相关的固有统计噪声的典型大小。

表~\ref{tab-chi} 以 \CTHERAII{} PDF 为例，总结了三个 LHC 喷注数据集在完成去关联并加入蒙特卡洛不确定度后 $\chi^2_E/N_{pt,E}$ 的缩减。例如对 ATLAS 7 TeV 喷注数据，去关联使 $\chi^2_E$ 降低约 90 单位以上，但仍需再加 MC 不确定度才能把 $\chi^2_E/N_{pt,E}$ 降到统计合理的水平（$N_{pt,E}=140$ 时从 1.68 到 1.31）。

\begin{table}[tb]
	\begin{tabular}{c|ccc}
		\hline
		\hline
        	&	\multicolumn{3}{c}{evaluated, {\bf CT14$_\mathrm{HERAII}$ NNLO}} 	\ \ \ \\
		\ \ \  $\chi^2_E/N_{pt,E}$ \ \ \	& \ \ \ original data \ \ \  &  \ \ \  $+$\,decorr. \ \ \    &  \ \ \  $+$\,0.5\% MC unc. \ \ \   \\
		\hline
		\ \ \ ATLAS, 7 TeV \ \ \   &  \ \ \  2.34 \ \ \   &  \ \ \  1.68 \ \ \   &  \ \ \  1.31 \ \ \   \\
		\ \ \ CMS, 7 TeV   \ \ \   &  \ \ \  1.58 \ \ \   &  \ \ \  1.45 \ \ \   &  \ \ \  1.35 \ \ \   \\
		\ \ \ CMS, 8 TeV   \ \ \   &  \ \ \  1.90 \ \ \   &  \ \ \  1.34 \ \ \   &  \ \ \  1.23 \ \ \   \\
		\hline
		\hline
	\end{tabular}
	\caption{CT18 实现的单举喷注产生数据的 $\chi^2_E/N_{pt,E}$，此处用 \CTHERAII{} PDF~\cite{Hou:2016nqm} 在几种误差处理情景下计算：首先是不实施任何去关联方案（``original data''）；其次是使用上文所述去关联方案（``$+$\, decorr.''）；最后是在去关联之上再添加 0.5\% 的整体非关联蒙特卡洛不确定度（``$+$\,0.5\% MC unc.''）。
	}
	\label{tab-chi}
\end{table}
